\begin{enumerate}

\item Moddiy nuqtaning garmonik tebranish tenglamasi $x = 0,04 \sin(4\pi t + \pi)$ (m) bo‘lsa, uning chastotasi (Hz) toping.
\begin{enumerate}[label=\Alph*)]
\item 2
\item 4
\item 6
\item 3
\end{enumerate}

\item Garmonik tebranishlar tenglamasi $x = 0,01 \cos 10t$ ([x]=m, [t]=s) dan foydalanib, bu tebranishlarning chastotasini aniqlang (Hz).
\begin{enumerate}[label=\Alph*)]
\item $5/\pi$
\item $10/\pi$
\item 5
\item 10
\end{enumerate}

\item Harakat tenglamasi $x = 0,06 \cos 100\pi t$ (m) ko‘rinishiga ega. Tebranish chastotasi qanday (Hz)?
\begin{enumerate}[label=\Alph*)]
\item 50
\item 4
\item 25
\item 20
\end{enumerate}

\item Garmonink tebranayotgan jismning harakat tenglamasi $x = 0,5 \cos 10\pi t$. Jism tebranishlarining chastotasi necha gers?
\begin{enumerate}[label=\Alph*)]
\item 10
\item 50
\item 0,5
\item 5
\end{enumerate}

\item Massasi 0,1 kg bo‘lgan prujinali mayatnikning tebranish qonuni $x = 0,5 \sin 10t$ (m) ko‘rinishga ega. Prujinaning bikrligi qanday (N/m)?
\begin{enumerate}[label=\Alph*)]
\item 2
\item 10
\item 5
\item 6,4
\end{enumerate}

\item Kosinus qonuni bo‘yicha garmonik tebranayotgan nuqtaning $\pi/3$ fazadagi siljishi 1 sm bo‘lsa, $5\pi/3$ fazadagi siljishi qanday bo‘ladi (sm)?
\begin{enumerate}[label=\Alph*)]
\item 0,5
\item 1
\item 1,25
\item 2,5
\end{enumerate}

\item Amplitudasi 1 mm, chastotasi 1000 Hz, boshlang‘ich fazasi $\pi/3$ bo‘lgan garmonik tebranishning tenglamasini tuzing (m).
\begin{enumerate}[label=\Alph*)]
\item $x = 0,001 \sin(2000\pi t)$
\item $x = \sin(1000\pi t + \pi/3)$
\item $x = \cos(1000\pi t + \pi/3)$
\item $x = 0,001 \cos(2000\pi t + \pi/3)$
\end{enumerate}

\item Amplitudasi 5 sm va davri 0,05 s bo‘lgan garmonik tebranishning tenglamasini tuzing (m).
\begin{enumerate}[label=\Alph*)]
\item $x = 0,5 \cos 40\pi t$
\item $x = 0,05 \sin 20\pi t$
\item $x = 0,05 \cos 20\pi t$
\item $x = 0,05 \sin 40\pi t$
\end{enumerate}

\item Garmonik tebranish fazasi ifodasini ko‘rsating.
\begin{enumerate}[label=\Alph*)]
\item $\cos(\omega t + a)$
\item $\omega t + a$
\item $\sin(\omega t + a)$
\item $\omega t$
\end{enumerate}

\item Keltirilganlar ichidan garmonik tebranishni xarakterlovchi ifodalarni toping: 1) $x = A \sin(\omega t + a)$; 2) $x = x_0 + v_0 t + at^2/2$; 3) $x = A \cos(\omega t + a)$; 4) $x = F/k$.
\begin{enumerate}[label=\Alph*)]
\item 1, 2
\item 1, 3
\item 2, 3
\item 3, 4
\end{enumerate}

\item $OX$ o‘qi bo‘ylab yuk tebranishi $x = 3 \cos(2t + \pi/2)$ tenglama bilan ifodalanadi. $t = 3$ s paytdagi tebranish fazasi aniqlansin.
\begin{enumerate}[label=\Alph*)]
\item 6
\item 3
\item $6 + \pi/2$
\item $\pi/2$
\end{enumerate}

\item Garmonik tebranish tenglamasi $x = 0,02 \cos \pi t$ ko‘rinishga ega, 0,25 s dan keyingi siljishdan 1,25 s dan keyingi siljish ayirmasini toping (sm).
\begin{enumerate}[label=\Alph*)]
\item 4,2
\item 1,4
\item 28
\item 2,8
\end{enumerate}

\item Keltirilgan grafikdan foydalanib, $t = 0,6$ s paytdagi tebranish fazasini aniqlang.
\begin{enumerate}[label=\Alph*)]
\item $0,6\pi$
\item $1,2\pi$
\item $0,3\pi$
\item $0,4\pi$
\end{enumerate}

\item Garmonik tebranishlarning amplitudasi 5 sm, boshlang‘ich fazasi $3\pi/2$ ga teng. Tebranishlar soni minutiga 120 marta tebranadi. Garmonik tebranishlarning tenglamasini ko‘rsating (m).
\begin{enumerate}[label=\Alph*)]
\item $x = 0,05 \cos(4\pi t + 3\pi/2)$
\item $x = 5 \cos(2\pi t + 3\pi/2)$
\item $x = 0,05 \sin 4\pi t$
\item $x = 5 \sin(4\pi t + \pi/2)$
\end{enumerate}

\item Jism $OX$ o‘qi bo‘ylab $x(t) = 6\pi \cos(\frac{\pi}{2}t + \frac{\pi}{4})$ (m) qonunga muvofiq tebranadi. Tebranishlar siklik chastotasini (rad/s) ko‘rsating.
\begin{enumerate}[label=\Alph*)]
\item $\pi/2$
\item $\pi/4$
\item $6\pi$
\item $\pi$
\end{enumerate}

\item Agar tebranishlar sinuslar qonuniga asosan sodir bo‘lsa, siljish qanday fazalarda modul bo‘yicha amplitudaning taqriban 0,87 qismiga teng bo‘ladi?
\begin{enumerate}[label=\Alph*)]
\item $\pi/3; 2\pi/3; 4\pi/3; 5\pi/3$ rad.
\item $\pi/4; \pi/2; 3\pi/4; 7\pi/4$ rad.
\item $\pi/6; \pi/3; \pi/2; 2\pi/3$ rad.
\item $\pi/2; \pi; \pi/2$ rad.
\end{enumerate}

\item Bir xil davrli va 5,0 sm, 2,0 sm amplitudali ikki garmonik tebranish bir to‘g‘ri chiziq bo‘yicha sodir bo‘lmoqda. Tebranish davri 1,2 s ga teng. Fazalar bir xil bo‘lganda tebranish boshlangandan 0,10 s keyingi natijaviy tebranish siljishini aniqlang (m).
\begin{enumerate}[label=\Alph*)]
\item $0,2 \cdot 10^{-2}$
\item $1,5 \cdot 10^{-3}$
\item $7 \cdot 10^{-2}$
\item $3,5 \cdot 10^{-2}$
\end{enumerate}

\item Tebranishlar amplitudasi 10 sm, chastotasi esa 0,5 Hz. Qancha vaqtdan keyin siljish 7,1 sm bo‘lishini aniqlang (Tebranish sinus qonuniga bo‘ysunadi).
\begin{enumerate}[label=\Alph*)]
\item 0,5 va 1,5 s
\item 0,25 va 2,4 s
\item 0,4 va 1,75 s
\item 0,25 va 1,75 s
\end{enumerate}

\item Tebranish davri 36 s bo‘lgan moddiy nuqta muvozanat vaziyatidan yarim amplitudagacha qancha vaqtda siljiydi (s)?
\begin{enumerate}[label=\Alph*)]
\item 9
\item 8
\item 4,5
\item 3
\end{enumerate}

\item Agar mayatnik 12 s davr bilan garmonik tebransa, chetki vaziyatidan amplitudaning yarmigacha bo‘lgan masofani necha sekundda bosib o‘tadi?
\begin{enumerate}[label=\Alph*)]
\item 2
\item 4
\item 6
\item 3
\end{enumerate}

\item Tebranayotgan jism amplituda yarmidan eng chekka nuqtagacha davrning qanday qismida siljiydi?
\begin{enumerate}[label=\Alph*)]
\item $T/12$
\item $T/3$
\item $T/4$
\item $T/6$
\end{enumerate}

\item Tebranayotgan jism muvozanat vaziyatidan amplitudagacha bo‘lgan masofani qancha vaqtda bosib o‘tadi?
\begin{enumerate}[label=\Alph*)]
\item $T/12$
\item $T/3$
\item $T/4$
\item $T/8$
\item $T/6$
\end{enumerate}

\item Tebranayotgan jism muvozanat vaziyatidan amplitudagacha borib yana muvozanat vaziyatiga qaytib kelishi uchun qancha vaqt ketadi (s)? Uning tebranish davri 6 s.
\begin{enumerate}[label=\Alph*)]
\item 3
\item 6
\item 2
\item 1
\end{enumerate}

\item Garmonik tebranayotgan jism muvozanat vaziyatidan chetki vaziyatgacha bo‘lgan masofaning birinchi yarmini davrning qanday qismida o‘tadi?
\begin{enumerate}[label=\Alph*)]
\item $1/6$
\item $1/4$
\item $1/12$
\item $1/8$
\end{enumerate}

\item Tebranayotgan jism tezligining tenglamasi $\vartheta = 5 \sin(3t + \pi/3)$ m/s ko‘rinishda bo‘lsa, jism siljishi tebranishlarining boshlang‘ich fazasi qanday?
\begin{enumerate}[label=\Alph*)]
\item $3t + \pi/3$
\item $3t$
\item 5
\item $-\pi/6$
\end{enumerate}

\item Tebranma harakat qilayotgan nuqtaning muvozanat vaziyatidan 4 sm uzoqlikdagi tezligi 6 sm/s, 3 sm uzoqlikdagi tezligi 8 sm/s ga teng. Tebranish amplitudasini toping (sm).
\begin{enumerate}[label=\Alph*)]
\item 4
\item 5
\item 6,7
\item 8
\end{enumerate}

\item Garmonik tebranayotgan nuqta tezligi vaqtga bog‘liq ravishda qanday o‘zgaradi?
\begin{enumerate}[label=\Alph*)]
\item o‘zgarmas saqlanadi.
\item vaqtga to‘g‘ri proporsional.
\item sinus yoki kosinus qonuni bo‘yicha.
\item vaqtga teskari proporsional.
\end{enumerate}

\item Mayatnikning tebranish qonuni $x = 0,1 \sin(\pi t + \pi)$ [m] ko‘rinishga ega. Tezlikning amplitudaviy qiymatini toping (m/s).
\begin{enumerate}[label=\Alph*)]
\item 0,1
\item $0,1\pi$
\item 1
\item $\pi$
\end{enumerate}

\item Moddiy nuqta koordinatasi $x = 2 + 3 \sin(4\pi t + \pi/2)$ qonun bo‘yicha o‘zgarmoqda. $t = 0$ vaqtda nuqta tezligi qanday bo‘lgan?
\begin{enumerate}[label=\Alph*)]
\item 5,3
\item 37,7
\item 12,0
\item 0
\end{enumerate}

\item $XOY$ tekislikda amplitudasi 10 mm bo‘lgan yassi garmonik to‘lqin $OY$ o‘qi bo‘ylab tarqalmoqda. Koordinatasi $y = 8$ mm bo‘lgan nuqtadagi zarraning muvozanat vaziyatidagi tezligi 50 sm/s. Shu zarraning muvozanat vaziyatidan siljishi 6 mm bo‘lgandagi tezligi (sm/s) qanday bo‘ladi?
\begin{enumerate}[label=\Alph*)]
\item 300
\item 500
\item 400
\item 600
\end{enumerate}

\item Garmonik tebranayotgan jismning siljishi moduli bo‘yicha eng katta qiymatga erishgan momentda tezligi o‘zining qanday qiymatiga erishadi?
\begin{enumerate}[label=\Alph*)]
\item eng katta
\item 0 ga teng bo‘ladi
\item o‘rta qiymatda bo‘ladi
\item aniqlab bo‘lmaydi
\end{enumerate}

\item Prujinaga osilgan 5 kg massali yuk 0,25 s davr bilan tebranmoqda va bir davrda 30 sm yo‘l bosmoqda. Yukning maksimal tezlanishini ($m/s^2$) toping.
\begin{enumerate}[label=\Alph*)]
\item 37,5
\item 47,3
\item 60
\item 9,6
\end{enumerate}

\item Matematik mayatnikning tebranish qonuni $x = 0,1 \sin(10\pi t + \pi)$ (m) ko‘rinishga ega. Tezlanishning amplitudaviy qiymatini toping ($m/s^2$).
\begin{enumerate}[label=\Alph*)]
\item $10\pi^2$
\item 100
\item 10
\item $\pi^2$
\end{enumerate}

\item Mayatnik $x = A \sin 100\pi t$ qonun bo‘yicha tebranmoqda. Uning tezlanish qonuni qanday ifodalanadi?
\begin{enumerate}[label=\Alph*)]
\item $a_x = 10 A \cos 100\pi t$
\item $a_y = -10^4 A \pi^2 t$
\item $a_x = 500 Ax$
\item $a_x = -10^4 \pi^2 x$
\end{enumerate}

\end{enumerate}



\begin{table}[!ht]
\centering
\begin{tabular}{|c|c|c|c|c|c|c|c|c|c|}
\hline
№ & J. & № & J. & № & J. & № & J. & № & J. \\ \hline
1 & A & 8 & D & 15 & A & 22 & C & 29 & D \\ \hline
2 & A & 9 & B & 16 & A & 23 & A & 30 & C \\ \hline
3 & A & 10 & B & 17 & D & 24 & C & 31 & B \\ \hline
4 & D & 11 & C & 18 & D & 25 & D & 32 & B \\ \hline
5 & B & 12 & D & 19 & D & 26 & B & 33 & A \\ \hline
6 & B & 13 & A & 20 & B & 27 & C & 34 & D \\ \hline
7 & D & 14 & A & 21 & D & 28 & B & & \\ \hline
\end{tabular}
\end{table}